\chapter{Completaciones estadísticas y campos sobre APP}
\label{chap:completaciones-cuanticas}

Las palabras Bose, Fermi, Pauli o fotón no designan por sí solas teoremas de
un sistema discreto. Para obtenerlos hay que especificar un espacio de estados,
una graduación, un producto y operadores. APP y TPK proporcionan un conjunto
de estados orientados y una paridad de hoja. A partir de esos datos pueden
construirse dos completaciones universales: el álgebra simétrica y el álgebra
exterior. Esta sección identifica exactamente lo que se sigue de la
construcción y separa la estadística algebraica de su interpretación física.

\section{Espacio lineal de estados asociado a una realización de partícula}

Sea \(S\) un conjunto finito de clases de rutas TPK a la profundidad elegida
y sea
\[
 \mathcal H_1=\ell^2(S)
\]
el espacio de Hilbert con base ortonormal \((e_s)_{s\in S}\). Una elección de
pesos positivos \(E_s\) define el operador diagonal
\[
 He_s=E_se_s.
\]
La elección de \(S\), de su cociente por simetrías y de los pesos forma parte
del modelo. Una vez fijada, las completaciones multipartícula son canónicas.

\section{Completaciones simétrica y exterior}

Definimos
\[
 \mathcal F_+(\mathcal H_1)=
 \bigoplus_{n\geq0}\operatorname{Sym}^n\mathcal H_1,
 \qquad
 \mathcal F_-(\mathcal H_1)=
 \bigoplus_{n\geq0}\bigwedge^n\mathcal H_1.
\]
La primera es la completación bosónica; la segunda, la fermiónica.

\begin{theorem}[Principio de exclusión de Pauli en la completación exterior]
Para todo modo \(s\in S\), los operadores de creación y aniquilación de la
completación exterior satisfacen
\[
 \{a_s,a_t^\dagger\}=\delta_{st}I,
 \qquad
 \{a_s,a_t\}=\{a_s^\dagger,a_t^\dagger\}=0.
\]
En particular, el operador de ocupación
\(N_s=a_s^\dagger a_s\) es un proyector:
\[
 N_s^2=N_s,
\]
y sus autovalores son \(0,1\).
\end{theorem}

\begin{proof}
La creación es el producto exterior por \(e_s\); la aniquilación es su
adjunto. La antisimetría implica \(e_s\wedge e_s=0\) y da las relaciones
canónicas por cálculo sobre la base exterior. Entonces
\[
N_s^2=a_s^\dagger a_sa_s^\dagger a_s
=a_s^\dagger(I-a_s^\dagger a_s)a_s=N_s.
\]
Todo proyector autoadjunto posee espectro contenido en \(\{0,1\}\).
\end{proof}

El principio de exclusión aparece así como teorema de la completación
exterior. La cuestión específicamente HMT es qué dato del transporte elige
esa completación para una clase de rutas.

\section{Graduación de hoja y álgebra de intercambio}

Supongamos que la hoja orientada define un carácter aditivo
\[
 p:\operatorname{Mor}(\mathsf P_{\rm TPK})\longrightarrow\mathbb Z/2\mathbb Z,
 \qquad p(\gamma\eta)=p(\gamma)+p(\eta).
\]
Sobre el espacio lineal graduado correspondiente se introduce la simetría
\[
 \tau(v\otimes w)=(-1)^{p(v)p(w)}w\otimes v.
\]

\begin{theorem}[Completación estadística determinada por la paridad]
El cociente del álgebra tensorial por las relaciones
\[
 v\otimes w-(-1)^{p(v)p(w)}w\otimes v
\]
es simétrico sobre el sector par y exterior sobre el sector impar. Dos estados
impares idénticos tienen producto nulo; los estados pares admiten ocupación
arbitraria.
\end{theorem}

\begin{proof}
Si \(p(v)=p(w)=0\), la relación es \(vw=wv\). Si
\(p(v)=p(w)=1\), es \(vw=-wv\); tomando \(v=w\) y trabajando en
característica cero se obtiene \(v^2=0\). Los sectores mixtos conmutan sin
signo adicional. Esta es la propiedad universal del álgebra simétrica
graduada.
\end{proof}

El teorema convierte una paridad composicional TPK en una regla exacta de
ocupación. Para obtener el teorema físico de espín--estadística hay que
construir además una teoría local lorentziana positiva y demostrar que la
paridad de hoja coincide con el espín bajo rotaciones. La igualdad no se
deduce únicamente de que ambos objetos tomen valores módulo dos.

\section{Holonomía de hoja y periodicidad de \texorpdfstring{\(2\pi/4\pi\)}{2pi/4pi}}

La paridad anterior admite una realización geométrica natural cuando una
extensión superviviente posee un retorno cerrado cuya cubierta de orientación
es no trivial. Sea \(\rho\) el generador del retorno y sea
\[
 \varepsilon:\langle\rho\rangle\longrightarrow\{-1,+1\}
\]
su carácter de holonomía, con \(\varepsilon(\rho)=-1\). El fibrado real
asociado es entonces la línea de orientación de una banda de Möbius: una
vuelta cierra la proyección geométrica, pero cambia la hoja levantada; dos
vueltas restauran la orientación.

Para escribir esta propiedad angularmente, fijemos una representación
unitaria \(U\) de la cubierta de rotaciones sobre el espacio graduado
\(\mathcal H_1=\mathcal H_{\bar0}\oplus\mathcal H_{\bar1}\), compatible con
la paridad de hoja en el sentido
\[
 U(2\pi)v=(-1)^{p(v)}v
\]
para todo vector homogéneo \(v\). La compatibilidad es una condición de la
realización: identifica el generador de la cubierta de orientación con el
carácter \(p\) ya construido sobre las rutas.

\begin{theorem}[Restauración de orientación y álgebra de intercambio]
\label{thm:dos-cuatro-pi-estadistica}
Bajo la compatibilidad precedente,
\[
 U(2\pi)|_{\mathcal H_{\bar0}}=I,
 \qquad
 U(2\pi)|_{\mathcal H_{\bar1}}=-I,
 \qquad
 U(4\pi)=I.
\]
Además, el mismo carácter \(p\) determina la simetría graduada
\[
 \tau(v\otimes w)=(-1)^{p(v)p(w)}w\otimes v.
\]
Por tanto, el sector de holonomía par se completa mediante potencias
simétricas y el sector de holonomía impar mediante potencias exteriores. En
el segundo se cumple la exclusión \(v\wedge v=0\).
\end{theorem}

\begin{proof}
La primera y la segunda igualdad son la condición de compatibilidad evaluada
en cada componente homogénea. Al iterar una vuelta,
\[
 U(4\pi)v=U(2\pi)^2v=(-1)^{2p(v)}v=v,
\]
lo que demuestra la restauración. La afirmación sobre el intercambio es el
teorema de completación por paridad de la sección anterior: para dos vectores
pares el signo es positivo, mientras que para dos impares es negativo.
\end{proof}

Conviene distinguir tres retornos. El cierre angular tras \(2\pi\) es un
retorno de la proyección; el cierre tras \(4\pi\) restaura la orientación
levantada; ninguno obliga por sí solo a que todas las coordenadas del registro
regresen a su valor inicial. Esta separación es la análoga dimensional de la
distinción HMT entre retorno de fase y retorno de estado.

\subsection{Las hojas suma y producto como soportes de las \(2\) completaciones}

La tabla APP contiene dos evaluaciones sobre un mismo soporte. La hoja
aditiva actúa por traslaciones y conserva la multiplicidad de las
continuaciones; en este sentido, acumula. La hoja multiplicativa es
permutativa sobre las unidades y contrae fibras sobre los no invertibles; en
este sentido, pliega. Esta incompatibilidad es anterior a la elección
estadística y explica por qué TPK debe conservar dos genealogías en lugar de
reducirlas a una sola palabra visible.

La asignación estadística no se obtiene, sin embargo, sustituyendo
verbalmente «suma» por «bosón» y «producto» por «fermión». El mapa exacto es
el carácter de paridad
\[
 p:\operatorname{Mor}(\mathsf P_{\rm TPK})\longrightarrow\mathbb Z/2\mathbb Z.
\]
Cuando el transporte de una extensión cae en el sector \(p=0\), su
completación admite ocupación acumulativa y es simétrica. Cuando cae en
\(p=1\), la reversión de hoja introduce el signo de intercambio y la
completación es exterior. La suma y el producto suministran las dos
geometrías locales; el registro, la orientación y el carácter \(p\) deciden en
qué sector se realiza una ruta.

Con un Hamiltoniano diagonal \(H\), energía \(\epsilon\), potencial químico
\(\mu\) y parámetro térmico \(\beta\), las ocupaciones medias de los dos
sectores son entonces
\[
 \bar n_{\rm par}(\epsilon)
 =\frac{1}{e^{\beta(\epsilon-\mu)}-1},
 \qquad
 \bar n_{\rm impar}(\epsilon)
 =\frac{1}{e^{\beta(\epsilon-\mu)}+1}.
\]
Las fórmulas de Bose--Einstein y Fermi--Dirac son, por tanto, las
evaluaciones termodinámicas de las completaciones simétrica y exterior
seleccionadas por el transporte. La periodicidad \(4\pi\) y la exclusión
fermiónica quedan coordinadas por el mismo carácter, pero ninguna de las dos
se infiere de una mera coincidencia numérica.

\section{Operadores de número y distribución de ocupación en la realización de Fock}

Si hay \(g\) modos y \(N\) partículas, las dimensiones de los sectores son
\[
 \Omega_F(g,N)=\binom gN,
 \qquad
 \Omega_B(g,N)=\binom{g+N-1}{N}.
\]
La primera fórmula cuenta subconjuntos de modos; la segunda, composiciones
débiles de \(N\) en \(g\) partes. Estos conteos son consecuencias exactas de
las dos completaciones y no requieren una hipótesis termodinámica.

\section{El complejo celular del toro APP}

El grafo de APP es el uno-esqueleto de una descomposición celular del toro
\(T^2\) con \(9\times9\) vértices. Sea
\[
 C^0\xrightarrow{d_0}C^1\xrightarrow{d_1}C^2
\]
su complejo de cocadenas reales. Elegido el producto escalar uniforme,
definimos la codiferencial \(\delta=d^*\) y el laplaciano
\[
 \Delta_k=d_{k-1}\delta_k+\delta_{k+1}d_k.
\]

\begin{theorem}[Sector armónico del toro digital]
La dimensión del espacio de uno-formas armónicas es dos:
\[
 \dim\ker\Delta_1=b_1(T^2)=2.
\]
Puede elegirse una base representada por los flujos uniformes horizontal y
vertical. Ambos son cerrados, cocerrados y no son gradientes de funciones
globales periódicas.
\end{theorem}

\begin{proof}
El teorema de Hodge para complejos celulares finitos identifica
\(\ker\Delta_1\) con \(H^1(T^2;\mathbb R)\). La cohomología del toro tiene
rango dos. Directamente, los flujos constantes horizontal y vertical tienen
divergencia y rotacional nulos; sus integrales sobre los dos ciclos
fundamentales son independientes, por lo que no son exactos.
\end{proof}

El espacio de conexiones módulo transformaciones de calibre es
\[
 C^1/\operatorname{im}d_0.
\]
En un complejo finito provisto del producto escalar anterior, cada órbita
admite un representante de Coulomb en \(\ker\delta_1\); dentro de la órbita,
la elección es única módulo el estabilizador \(\ker d_0\) del parámetro de
calibre. El cociente cohomológico correcto es
\[
 \ker d_1/\operatorname{im}d_0
 \simeq \ker\Delta_1 .
\]
Sus dos clases armónicas constituyen un candidato matemático natural para
dos modos globales libres. No son todavía las dos polarizaciones locales de
un fotón en espacio-tiempo \(3+1\): para esa identificación se necesita un
límite dinámico, una relación de dispersión y un grupo físico de calibre.

\section{Ley de Gauss discreta}

Para una 1-cocadena orientada \(F\) y un conjunto de vértices \(V\), la
divergencia se define por
\[
 (\operatorname{div}F)(v)=
 \sum_{e:v\to w}F(e)-\sum_{e:w\to v}F(e).
\]
Al sumar sobre \(V\), toda arista interior aparece una vez con cada signo y
se cancela. Por tanto,
\[
 \boxed{
 \sum_{v\in V}\operatorname{div}F(v)
 =\operatorname{Flux}_{\partial V}(F).}
\]
Esta identidad es una ley de Gauss exacta en el grafo. Interpretar
\(\operatorname{div}F\) como carga eléctrica exige el carácter de carga y la
realización dimensional construidos en Mecánica Dimensional.

\section{Aparato de medición reversible}

Sea \(S\) un conjunto de estados y sea \(r:S\to\mathbb Z/9\mathbb Z\) una
lectura. El acoplamiento a un registro nonádico se define por
\[
 U_r:S\times\mathbb Z/9\mathbb Z\longrightarrow
 S\times\mathbb Z/9\mathbb Z,
 \qquad
 U_r(s,a)=(s,a+r(s)).
\]

\begin{proposition}[Reversibilidad del acoplamiento]
La aplicación \(U_r\) es una permutación y
\[
 U_r^{-1}(s,a)=(s,a-r(s)).
\]
La pérdida de información sólo aparece después de aplicar un lector al
registro o de condicionar el estado por el resultado.
\end{proposition}

\begin{proof}
La fórmula propuesta para la inversa satisface ambas composiciones por la
ley de grupo de \(\mathbb Z/9\mathbb Z\).
\end{proof}

Esta proposición formaliza la definición operativa de medición: acoplar,
leer y restringir futuros compatibles son tres aplicaciones diferentes. El
acoplamiento completo es reversible; el cociente observacional puede no serlo.

\section{Correspondencias algebraicas y realizaciones físicas}

Las completaciones de Fock, la exclusión condicionada a paridad, el sector
armónico de APP, Gauss discreto y la reversibilidad del aparato son teoremas
matemáticos completos. Sus identificaciones físicas forman un diagrama que
debe conmutar:
\[
 \text{paridad de hoja}
 \longrightarrow\text{intercambio graduado}
 \longrightarrow\text{CAR/CCR}
 \longrightarrow\text{estadística física},
\]
\[
 H^1(\mathcal A_9)
 \longrightarrow\text{campo discreto de calibre}
 \longrightarrow\text{modo sin masa}
 \longrightarrow\text{fotón}.
\]
Las dos primeras flechas de cada línea poseen construcciones explícitas; las
últimas requieren la dinámica y el límite físico correspondientes. Esta
separación permite desarrollar los puentes sin confundir una analogía de
cardinalidad con una representación.
